% !PRINTED 493
\weberpage{493}

We prove further the theorem:

5. \gesperrt{Every algebraic number $\Theta$ can be transformed by multiplication
with a natural number into an integral algebraic number.}

For if
\[
\Theta^m + A_1 \Theta^{m-1} + \cdots + A_m = 0,
\]
and $A_1, \ldots, A_m$ are rational numbers with the common denominator $a$, then
by multiplication with $a^m$ one obtains:
\[
(a\Theta)^m + A_1 a (a\Theta)^{m-1} + A_2 a^2 (a\Theta)^{m-2} + \cdots + A_m a^m = 0,
\]
from which it follows that $a\Theta$ is an \gesperrt{integer}.

\begin{center}
\S.\,134.

\gesperrt{Algebraic Fields.}
\end{center}

In the thirteenth chapter of the first volume we have seen how, from every
root $\Theta$ of an equation $f(\Theta) = 0$ of degree $n$ irreducible in any
field $\Omega$, one derives an algebraic field $\Omega(\Theta)$ over $\Omega$.
The $n$ fields derived from the $n$ roots of this equation, which may also in
part or all be identical, we have called \gesperrt{conjugate fields}.

If we denote by $R$ the field of rational numbers, then by our definition every
algebraic number of degree $n$, $\Theta$, gives rise to an algebraic field
$R(\Theta)$ over $R$, which from now on we shall call briefly an
algebraic number field of degree $n$. We have also already proved
earlier (Bd.\,I, \S.\,143) that one can always determine an \gesperrt{algebraic number
field which contains a finite number of arbitrarily given algebraic numbers.}

This theorem is emphasised at this point in order to point out that the
generality of a consideration concerning any finite number of algebraic numbers
is not impaired by supposing these numbers all to lie in one algebraic number
field.

Every number $\omega$ of such a field can be represented as an integral
function $\varphi(\Theta)$ of $\Theta$ with rational coefficients, and to every
number $\omega$ there corresponds in each of the $n$ conjugate fields a
determinate number. These \gesperrt{conjugate numbers} can in part be equal to one
another, and we have accordingly primitive and
%CONT>>
